\scanpage{21}
\uline{Lemme} 17 Soit $f:X\longrightarrow S$ propre, \inserted{$S$ noeth.,} et soit $\underline E$ coh. sur $X$. Alors il existe un ouvert dense $U$ de $S$ tel que $\forall\,s\in S$, et $\underline f=f_1,\ldots f_{\unsure{n|r}}$ suite $\mathcal O_{S,s}$-régulière, l'hom naturel
\[
 R^jf_*(E)_s\otimes\mathcal O_{S,s}/\underline f\mathcal O_{S,s}
 \longrightarrow R^j f^{(s)}_*(E^{(s)}\otimes\deleted{\unsure{A}}\mathcal O_{S,s}/\underline f\mathcal O_{S,s})
 \tag*{$\alpha$)}
\]
soit un isom pour tt $j$.
\margin{($f^{(s)}:X^{(s)}=X\times_S\operatorname{Spec}\mathcal O_{S,s}\longrightarrow\operatorname{Spec}\mathcal O_{S,s}$, $E^{(s)}$ induit par $E$ sur $X^{(s)}$)}

\uline{Dém} Utilisant le lemme 16, on voit qu'on peut supposer déjà que
\[
 E^{(s)}\otimes\mathcal O_{S,s}/\underline f\mathcal O_{S,s}
 \overset{\sim}{\longleftarrow}E^{(s)}\overset{\mathbb L}{\otimes}\underline{O}_{S,s}/\underline f\mathcal O_{S,s}.
\]
De \unsure{m\textsuperscript{e}|même}, on peut supposer \deleted{que} (faisant $X=S$, \unsure{..} et représentant $Rf_*(E)$ par un complexe borné à composantes cohérentes $K^\bullet$ sur $S$) que
\[
 K^\bullet\otimes \underline{O}_{S,s}/\underline f\mathcal O_{S,s}
 =K^\bullet\overset{\mathbb L}{\otimes}\mathcal O_{S,s}/\underline f\mathcal O_{S,s}
\]
de sorte que l'hom $\alpha$, grâce à l'isom de projection (où $M\unsure{:|=}A_p/\underline f A_p$)
\[
 Rf_*(\underbrace{E\overset{\mathbb L}{\otimes}_{\mathcal O_S}M}_{E\otimes_{\mathcal O_S}M})
 \longrightarrow\underbrace{Rf_*(E)}_{K^\bullet}\overset{\mathbb L}{\otimes}_{\underline{O}_S}M,
\]
s'identifie à
\[
 H^j(K^\bullet)\otimes M\unsure{-|\longrightarrow} H^j(\overbrace{K^\bullet\otimes M}^{\unsure{\prime\prime|=}K^\bullet\overset{\mathbb L}{\otimes}M}).
\]
Si on \unsure{rapetisse|restreint} encore $U$, on peut supposer
\[
 \operatorname{Tor}_i^{\mathcal O_S}(H^j(K^\bullet),M)=0\qquad\text{pour }i>0,\ \text{tout }j,
\]
et on sait qu'alors l'hom précédent est un isom pour tt $j$. \quad cqfd.
% archive sheet number: 20
